\documentclass{standalone}
\usepackage{pgfplots}
\usepackage{tikz-3dplot}
\usetikzlibrary{fadings}
\pgfplotsset{compat = newest}

%Zmienne do bardziej intuicyjnej modyfikacji widoku trójwymiarowego
\def\upDown{80}
\def\leftRight{-35}

%Definicje rozmiarów belek
\def\dlugosc{45}
\def\szerokosc{4}

%Wymiary obrazka
\def\XO{-17}
\def\XI{7}

%Definicje kolorów
\definecolor{BEIGE}{RGB}{249,242,231}
\definecolor{BEIGEDARK}{RGB}{215,162,77}
\definecolor{GRAYLIGHT}{RGB}{229,229,231}
\definecolor{BLUEDARK}{RGB}{43,48,93}
\definecolor{BLUE}{RGB}{65,91,173}

\begin{document}
    \begin{tikzpicture}
        %Tło fade z lewej na prawą (ciemnoniebieski do jasnoniebieskiego), z dołu do góry (ciemnoniebieski do alpha)
        \fill[shading = axis, left color = BLUE!75!GRAYLIGHT!100, right color = BLUEDARK!80!black!100] (\XO,-3.5) rectangle (\XI,6);
        \fill[BLUEDARK, path fading = north] (\XO,-4) rectangle (\XI,6);
        %Płaszczyzna pozioma
        \fill[BLUE] (\XO,-6) rectangle (\XI,-3.5);
        \fill[GRAYLIGHT!20!BLUEDARK!100, path fading = west] (\XO,-6) rectangle (\XI,-3.5);
        \fill[BLUEDARK!80!black!100, path fading = south] (\XO,-6) rectangle (\XI,-3.5);

        %Cień kuli
        \tikzfading[name = fade out, inner color = transparent!50, outer color = transparent!100];
        \filldraw[black, path fading = fade out] (-10,-5) rectangle (10,-3.8);

        %Cienie liter
        \tikzfading[name = fade out, inner color = transparent!50, outer color = transparent!100];
        \filldraw[black, path fading = fade out] (-20,-5.75) rectangle (-4,-3.25);

        %Wnętrze kuli
        \path[fill = BLUEDARK] (0,0) circle (4.35);

        %Refleksy świetlne w kuli
        \tikzfading[name = fade out, inner color = transparent!25, outer color = transparent!100];
        \filldraw [BLUE!95!white!95,path fading = fade out] (0,-4) rectangle (4,1);
        \tikzfading[name = fade out, inner color = transparent!25, outer color = transparent!100];
        \filldraw [BLUE!95!white!95,path fading = fade out] (-4,-1) rectangle (0,4);

        %Obracająca się obręcz
        \begin{scope}
            %Pomocniczy okrąg
            \node[circle, minimum size = 5.5cm] (N) at (0,0){};
            % 1/6 okręgu - sztuczny obrót obręczy o 90 stopni
            \foreach \i in {0, 5, ..., 60}{
                \begin{scope}[rotate = \i]
                    \coordinate (A) at (N.\i);
                    \coordinate (B) at ([xshift = \dlugosc] N.\i);
                    \coordinate (C) at ([xshift = \dlugosc, yshift = \szerokosc] N.\i);
                    \coordinate (D) at ([yshift = \szerokosc] N.\i);
                    \coordinate (E) at ([xshift = \dlugosc - \dlugosc * (\i / 60)] N.\i);
                    \coordinate (F) at ([xshift = \dlugosc - \dlugosc * (\i / 60),  yshift = \szerokosc] N.\i);
                    \path[shading = axis, top color = BEIGEDARK!100!black!125, bottom color = BEIGE, shading angle = \i + 100] (A) -- (E) -- (F) -- (D) -- cycle;
                    \path[shading = axis, top color = BEIGEDARK!100!black!125, bottom color = BEIGE, shading angle = \i + 100] (E) -- (F) -- (C) -- (B) -- cycle;
                \end{scope}
            }
            % 2/6 okręgu
            \foreach \i in {65, 70, ..., 120}{
                \begin{scope}[rotate = \i]
                    \coordinate (A) at (N.\i);
                    \coordinate (B) at ([xshift = \dlugosc] N.\i);
                    \coordinate (C) at ([xshift = \dlugosc, yshift = \szerokosc] N.\i);
                    \coordinate (D) at ([yshift = \szerokosc] N.\i);
                    \coordinate (E) at ([xshift = \dlugosc - \dlugosc * ((\i - 60) / 60)] N.\i);
                    \coordinate (F) at ([xshift = \dlugosc - \dlugosc * ((\i - 60) / 60), yshift = \szerokosc] N.\i);
                    \path[shading = axis, top color = BEIGEDARK!100!black!125, bottom color = BEIGE, shading angle = \i + 100] (A) -- (E) -- (F) -- (D) -- cycle;
                    \path[shading = axis, top color = BEIGEDARK!100!black!125, bottom color = BEIGE, shading angle = \i + 100] (E) -- (F) -- (C) -- (B) -- cycle;
                \end{scope}
            }
            % 3/6 okręgu
            \foreach \i in {125, 130, ..., 180}{
                \begin{scope}[rotate = \i]
                    \coordinate (A) at (N.\i);
                    \coordinate (B) at ([xshift = \dlugosc] N.\i);
                    \coordinate (C) at ([xshift = \dlugosc, yshift = \szerokosc] N.\i);
                    \coordinate (D) at ([yshift = \szerokosc] N.\i);
                    \coordinate (E) at ([xshift = \dlugosc - \dlugosc * ((\i - 120) / 60)] N.\i);
                    \coordinate (F) at ([xshift = \dlugosc - \dlugosc * ((\i - 120) / 60), yshift = \szerokosc] N.\i);
                    \path[shading = axis, top color = BEIGEDARK!100!black!125, bottom color = BEIGE, shading angle = \i + 100] (A) -- (E) -- (F) -- (D) -- cycle;
                    \path[shading = axis, top color = BEIGEDARK!100!black!125, bottom color = BEIGE, shading angle = \i + 100] (E) -- (F) -- (C) -- (B) -- cycle;
                \end{scope}
            }
            % 4/6 okręgu
            \foreach \i in {185, 190, ..., 240}{
                \begin{scope}[rotate = \i]
                    \coordinate (A) at (N.\i);
                    \coordinate (B) at ([xshift = \dlugosc] N.\i);
                    \coordinate (C) at ([xshift = \dlugosc, yshift = \szerokosc] N.\i);
                    \coordinate (D) at ([yshift = \szerokosc] N.\i);
                    \coordinate (E) at ([xshift = \dlugosc - \dlugosc * ((\i - 180) / 60)] N.\i);
                    \coordinate (F) at ([xshift = \dlugosc - \dlugosc * ((\i - 180) / 60), yshift = \szerokosc] N.\i);
                    \path[shading = axis, top color = BEIGEDARK!100!black!125, bottom color = BEIGE, shading angle = \i + 100] (A) -- (E) -- (F) -- (D) -- cycle;
                    \path[shading = axis, top color = BEIGEDARK!100!black!125, bottom color = BEIGE, shading angle = \i + 100] (E) -- (F) -- (C) -- (B) -- cycle;
                 \end{scope}
            }
            % 5/6 okręgu
            \foreach \i in {245, 250, ..., 300}{
                \begin{scope}[rotate = \i]
                    \coordinate (A) at (N.\i);
                    \coordinate (B) at ([xshift = \dlugosc] N.\i);
                    \coordinate (C) at ([xshift = \dlugosc, yshift = \szerokosc] N.\i);
                    \coordinate (D) at ([yshift = \szerokosc] N.\i);
                    \coordinate (E) at ([xshift = \dlugosc - \dlugosc * ((\i - 240) / 60)] N.\i);
                    \coordinate (F) at ([xshift = \dlugosc - \dlugosc * ((\i - 240) / 60), yshift = \szerokosc] N.\i);
                    \path[shading = axis, top color = BEIGEDARK!100!black!125, bottom color = BEIGE, shading angle = \i + 100] (A) -- (E) -- (F) -- (D) -- cycle;
                    \path[shading = axis, top color = BEIGEDARK!100!black!125, bottom color = BEIGE, shading angle = \i + 100] (E) -- (F) -- (C) -- (B) -- cycle;
                \end{scope}
            }
            % 6/6 okręgu
            \foreach \i in {305, 310, ..., 355}{
                \begin{scope}[rotate = \i]
                    \coordinate (A) at (N.\i);
                    \coordinate (B) at ([xshift = \dlugosc] N.\i);
                    \coordinate (C) at ([xshift = \dlugosc, yshift = \szerokosc] N.\i);
                    \coordinate (D) at ([yshift = \szerokosc] N.\i);
                    \coordinate (E) at ([xshift = \dlugosc - \dlugosc * ((\i - 300) / 60)] N.\i);
                    \coordinate (F) at ([xshift = \dlugosc - \dlugosc * ((\i - 300) / 60), yshift = \szerokosc] N.\i);
                    \path[shading = axis, top color = BEIGEDARK!100!black!125, bottom color = BEIGE, shading angle = \i + 100] (A) -- (E) -- (F) -- (D) -- cycle;
                    \path[shading = axis, top color = BEIGEDARK!100!black!125, bottom color = BEIGE, shading angle = \i + 100] (E) -- (F) -- (C) -- (B) -- cycle;
                \end{scope}
            }
        \end{scope}

        % Cień w środku kuli
        \tikzfading[name = fade out, inner color = transparent!25, outer color = transparent!100];
        \filldraw [black, path fading = fade out] (-3.75,-3) rectangle (-2.5,2);
        \tikzfading[name = fade out, inner color = transparent!60, outer color = transparent!100];
        \filldraw [black, path fading = fade out] (-3.75,-4) rectangle (2,0);

        % Pęknięta kula
        \path[shading = ball, ball color = BLUE] (120:4.36) arc (120:320:4.36) arc (-50:-120:3) arc (-80:-200:2) arc (-120: -200: 2) arc (-180: -200: 1) -- cycle;

        % Wymiar obrazka
        \path (\XO,-6) rectangle (\XI,6);

        % Trójwymiarowa część obrazka
        \begin{scope}[shift = {(-7.5,0,9)}, scale=0.6]
            \tdplotsetmaincoords{\upDown}{\leftRight}
            \coordinate (Shift) at (0,0,0);
            \tdplotsetrotatedcoordsorigin{(Shift)};
            \tdplotsetrotatedcoords{-20}{0}{0}
            \begin{scope}[tdplot_main_coords, tdplot_rotated_coords]
                %Napis TikZ
                \path[shading = axis, top color = BLUEDARK!90!BLUE!100, bottom color = BLUEDARK!90!BLUE!100, middle color = BLUE!100!black!85, shading angle = -45]
                (-6,1,-2) -- (-6,1,4) -- (-7.5,1,4) -- (-7.5,1,5.25) -- (-3.5,1,5.25) -- (-3.5,1,4) -- (-4.85,1,4) -- (-4.85,1,-2) -- cycle

                (-2.75,1,-2) -- (-2.75,1,3) -- (-1.55,1,3) -- (-1.55,1,-2) -- cycle

                (-2.75,1,4) -- (-2.75,1,5.15) -- (-1.55,1,5.15) -- (-1.55,1,4) -- cycle

                (-0.5,1,-2) -- (-0.5,1,5.15) -- (0.75,1,5.15) -- (0.75,1,0.85) -- (2.3,1,3.15) -- (3.5,1,3.15) -- (2,1,1.15) -- (3.85,1,-2) -- (2.5,1,-2) -- (1.45,1,0.25) -- (0.75,1,-0.75) -- (0.75,1,-2) -- cycle

                (4.5,1,-2) -- (4.5,1,-1) -- (7,1,4.15) -- (4.4,1,4.15) -- (4.4,1,5.25) -- (8,1,5.25) -- (8,1,4.15) -- (5.45,1,-1) -- (8,1,-1) -- (8,1,-2) -- cycle;

                %(-6,0,-2) -- (-4.85,0,-2) -- (-4.85,1,-2) -- (-6,1,-2) -- cycle - podstawa 'T'
                %(4.5,1,-2) -- (4.5,0,-2) -- (8,0,-2) -- (8,1,-2) -- cycle - podstawa 'Z'
                %Przy wspólnym malowaniu ścian ^^^ powodują błędy. I tak są niewidoczne, zatem można się ich pozbyć
                \path[shading = axis, top color = BLUEDARK!90!BLUE!110, bottom color = BLUEDARK!90!BLUE!100, shading angle = 5, BLUEDARK!90!BLUE!110]
                (-6,1,4) -- (-6,0,4) -- (-7.5,0,4) -- (-7.5,1,4) -- cycle
                (-7.5,0,4) -- (-7.5,1,4) -- (-7.5,1,5.25) -- (-7.5,0,5.25) -- cycle
                (-7.5,1,5.25) -- (-7.5,0,5.25) -- (-3.5,0,5.25) -- (-3.5,1,5.25) -- cycle
                (-3.5,0,5.25) -- (-3.5,1,5.25) -- (-3.5,1,4) -- (-3.5,0,4) -- cycle
                (-3.5,1,4) -- (-3.5,0,4) -- (-4.85,0,4) -- (-4.85,1,4) -- cycle
                (-4.85,0,4) -- (-4.85,1,4) -- (-4.85,1,-2) -- (-4.85,0,-2) -- cycle
                (-6,0,-2) -- (-6,1,-2) -- (-6,1,4) -- (-6,0,4) -- cycle

                (-2.75,0,-2) -- (-1.55,0,-2) -- (-1.55,1,-2) -- (-2.75,1,-2) -- cycle
                (-1.55,0,-2) -- (-1.55,1,-2) -- (-1.55,1,3) -- (-1.55,0,3) -- cycle
                (-1.55,1,3) -- (-1.55,0,3) -- (-2.75,0,3) -- (-2.75,1,3) -- cycle
                (-2.75,0,3) -- (-2.75,1,3) -- (-2.75,1,-2) -- (-2.75,0,-2) -- cycle
                (-2.75,0,4) -- (-1.55,0,4) -- (-1.55,1,4) -- (-2.75,1,4) -- cycle
                (-1.55,0,4) -- (-1.55,1,4) -- (-1.55,1,5.15) -- (-1.55,0,5.15) -- cycle
                (-1.55,1,5.15) -- (-1.55,0,5.15) -- (-2.75,0,5.15) -- (-2.75,1,5.15) -- cycle
                (-2.75,0,5.15) -- (-2.75,1,5.15) -- (-2.75,1,4) -- (-2.75,0,4) -- cycle
                (-0.5,0,-2) -- (0.75,0,-2) -- (0.75,1,-2) -- (-0.5,1,-2) -- cycle
                (0.75,0,-2) -- (0.75,1,-2) -- (0.75,1,-0.75) -- (0.75,0,-0.75) -- cycle
                (0.75,1,-0.75) -- (0.75,0,-0.75) -- (1.45,0,0.25) -- (1.45,1,0.25) -- cycle
                (1.45,0,0.25) -- (1.45,1,0.25) -- (2.5,1,-2) -- (2.5,0,-2) -- cycle
                (2.5,1,-2) -- (2.5,0,-2) -- (3.85,0,-2) -- (3.85,1,-2) -- cycle
                (3.85,0,-2) -- (3.85,1,-2) -- (2,1,1.15) -- (2,0,1.15) -- cycle
                (2,1,1.15) -- (2,0,1.15) -- (3.5,0,3.15) -- (3.5,1,3.15) -- cycle
                (3.5,0,3.15) -- (3.5,1,3.15) -- (2.3,1,3.15) -- (2.3,0,3.15) -- cycle
                (2.3,1,3.15) -- (2.3,0,3.15) -- (0.75,0,0.85) -- (0.75,1,0.85) -- cycle
                (0.75,0,0.85) -- (0.75,1,0.85) -- (0.75,1,5.15) -- (0.75,0,5.15) -- cycle
                (0.75,1,5.15) -- (0.75,0,5.15) -- (-0.5,0,5.15) -- (-0.5,1,5.15) -- cycle
                (-0.5,0,5.15) -- (-0.5,1,5.15) -- (-0.5,1,-2) -- (-0.5,0,-2) -- cycle

                (4.5,1,-2) -- (4.5,1,-1) -- (4.5,0,-1) -- (4.5,0,-2) -- cycle

                (8,0,-2) -- (8,1,-2) -- (8,1,-1) -- (8,0,-1) -- cycle
                (8,1,-1) -- (8,0,-1) -- (5.45,0,-1) -- (5.45,1,-1) -- cycle
                (5.45,0,-1) -- (5.45,1,-1) -- (8,1,4.15) -- (8,0,4.15) -- cycle
                (8,1,4.15) -- (8,0,4.15) -- (8,0,5.25) -- (8,1,5.25) -- cycle
                (8,0,5.25) -- (8,1,5.25) -- (4.4,1,5.25) -- (4.4,0,5.25) -- cycle
                (4.4,1,5.25) -- (4.4,0,5.25) -- (4.4,0,4.15) -- (4.4,1,4.15) -- cycle
                (4.4,0,4.15) -- (4.4,1,4.15) -- (7,1,4.15) -- (7,0,4.15) -- cycle
                (7,1,4.15) -- (7,0,4.15) -- (4.5,0,-1) -- (4.5,1,-1) -- cycle;

                \path[shading = axis, top color = BLUEDARK!90!BLUE!90, bottom color = BLUEDARK!90!BLUE!90, middle color = BLUE!100!black!65, shading angle = 30]
                (-6,0,-2) -- (-6,0,4) -- (-7.5,0,4) -- (-7.5,0,5.25) -- (-3.5,0,5.25) -- (-3.5,0,4) -- (-4.85,0,4) -- (-4.85,0,-2) -- cycle

                (-2.75,0,-2) -- (-2.75,0,3) -- (-1.55,0,3) -- (-1.55,0,-2) -- cycle

                (-2.75,0,4) -- (-2.75,0,5.15) -- (-1.55,0,5.15) -- (-1.55,0,4) -- cycle

                (-0.5,0,-2) -- (-0.5,0,5.15) -- (0.75,0,5.15) -- (0.75,0,0.85) -- (2.3,0,3.15) -- (3.5,0,3.15) -- (2,0,1.15) -- (3.85,0,-2) -- (2.5,0,-2) -- (1.45,0,0.25) -- (0.75,0,-0.75) -- (0.75,0,-2) -- cycle

                (4.5,0,-2) -- (4.5,0,-1) -- (7,0,4.15) -- (4.4,0,4.15) -- (4.4,0,5.25) -- (8,0,5.25) -- (8,0,4.15) -- (5.45,0,-1) -- (8,0,-1) -- (8,0,-2) -- cycle;
            \end{scope}
        \end{scope}
    \end{tikzpicture}
\end{document}
